\chapter{Lean 4 Formal Verification}

\section{Why Formal Verification?}

LaTeX templates can harbour subtle bugs that escape manual review: wrong
margin calculations, inconsistent doctype-to-class mappings, missing
translation keys, or broken fallback chains.  These errors surface only
when a user hits the specific combination of options that triggers them.

OmniLaTeX addresses this with Lean~4 formal verification.  Properties of
the template system are expressed as mathematical theorems and proven
correct at compile time.  If a pull request breaks a mapping or
introduces a contradiction, \texttt{lake build} fails before the LaTeX
code is ever compiled.

The proofs live in \texttt{specs/proofs/} and are verified by CI on every
push to \texttt{main}/\texttt{master} (see Chapter~\ref{ch:ci-cd}).

\section{Proof Modules}

All ten proof modules are imported through the root file
\texttt{specs/proofs/OmniLaTeXProofs.lean}.  Each module covers a
distinct aspect of the template system:

\begin{table}[htbp]
  \centering
  \begin{tabular}{@{} l r p{6.5cm} @{}}
    \toprule
    \textbf{Module} & \textbf{Theorems} & \textbf{What It Proves} \\
    \midrule
    BuildModes & 5 &
      Build mode properties: ultra skips bibliography,
      prod validates datamodel, dev has enough passes,
      all modes enable shell-escape, strictness ordering. \\[3pt]
    BuildSystem & 9 &
      Cache correctness, 43-example count, parallelism bounds
      ($0 < \text{workers} \leq \text{examples}$), production is
      default. \\[3pt]
    DoctypeClassMapping & 7 &
      Deterministic doctype$\to$class mapping, 60 aliases across
      3 KOMA classes (34 scrartcl, 6 scrbook, 20 scrreprt),
      covering theorem, scrartcl is largest. \\[3pt]
    DocumentSettings & 8 &
      KOMA class partition: 15 article-class, 6 report-class,
      5 book-class doctypes sum to 26, each non-empty. \\[3pt]
    DoctypeResolution & 4 &
      Alias resolution is deterministic and total over 55 known
      aliases; profile-to-class consistency. \\[3pt]
    FloatInvariant & 2 &
      Float placement never exceeds its section boundary;
      bounded deferral invariant. \\[3pt]
    FontHierarchy & 4 &
      Font sizes form a strict total order: irreflexive, asymmetric,
      transitive, and connex over 9 sizes. \\[3pt]
    I18nCompleteness & 3 &
      Translation key count is constant across all 18 languages;
      $18 \times 47 = 846$ total keys. \\[3pt]
    LanguageFallback & 7 &
      18 primary + 26 secondary = 44 languages; all fallbacks
      terminate at English; fallback is idempotent. \\[3pt]
    PageGeometry & 5 &
      Horizontal and vertical page balance equations hold for
      oneside and twoside layouts; DIV textwidth formula;
      caption width bounds. \\
    \bottomrule
  \end{tabular}
  \caption{Lean~4 proof modules and their properties}
\end{table}

\section{Proof Structure}

Lean~4 proofs use a small set of foundational constructs:

\textbf{Inductive types} define the domain.  For example, the three build
modes are declared as:

\begin{minted}{lean}
inductive BuildMode where
  | dev : BuildMode
  | prod : BuildMode
  | ultra : BuildMode
  deriving DecidableEq, Repr
\end{minted}

\textbf{Structures} bundle related fields:

\begin{minted}{lean}
structure BuildConfig where
  mode : BuildMode
  maxPasses : Nat
  runBiber : Bool
  deriving Repr
\end{minted}

\textbf{Function definitions} provide the mapping logic:

\begin{minted}{lean}
def buildConfigFor : BuildMode → BuildConfig
  | .dev => ⟨.dev, 10, true, true, false, true, false⟩
  | .prod => ⟨.prod, 5, true, true, true, true, true⟩
  | .ultra => ⟨.ultra, 1, false, false, false, true, false⟩
\end{minted}

\textbf{Theorems} state the properties to prove.  The proof follows after
\texttt{:= by}:

\begin{minted}{lean}
theorem ultra_no_bib :
    (buildConfigFor .ultra).runBiber = false ∧
    (buildConfigFor .ultra).runBib2gls = false := by
  simp [buildConfigFor]
\end{minted}

The hierarchy within each module follows the convention
\texttt{theorem} for main results, with \texttt{lemma} and
\texttt{corollary} used for derived statements.

\section{Reading Lean Proofs}

Key patterns you will encounter in the proof files:

\begin{description}
  \item[{\texttt{def}}] Function or constant definitions.
  \item[{\texttt{theorem} / \texttt{lemma}}] Statements to be proven.
  \item[{\texttt{by}}] Opens a tactic-mode proof block.
  \item[{\texttt{simp [name]}}] Rewrites using the definition of
        \texttt{name}.
  \item[{\texttt{decide}}] Automatically proves decidable propositions
        (finite enumeration, equality checks).
  \item[{\texttt{omega}}] Solves linear arithmetic over integers or
        naturals.
  \item[{\texttt{cases}}] Pattern-matches on an inductive type.
  \item[{\texttt{intro} / \texttt{exact}}] Introduces hypotheses and
        provides exact proof terms.
  \item[{\texttt{have}}] Introduces a local lemma.
  \item[{\texttt{rw} / \texttt{simp\_all}}] Rewrites equations in the
        goal or all hypotheses.
\end{description}

\section{Zero sorry Policy}

All theorems across all ten modules are fully proven.  No \texttt{sorry}
or \texttt{admit} appears anywhere in the proof source.  A \texttt{sorry}
in Lean~4 is a placeholder that makes the type checker accept an
unproven theorem -- essentially an assertion without evidence.

The CI workflow (\texttt{lean4-ci.yml}) enforces this by running
\texttt{grep -rq 'sorry' specs/proofs/} after every build.  If any
\texttt{sorry} is found, the step summary emits a warning, making
incomplete proofs visible in the Actions UI.

\section{Extending Proofs}

To add a new proof module:

\begin{enumerate}
  \item Create a new \texttt{.lean} file in
        \texttt{specs/proofs/OmniLaTeXProofs/}.
  \item Define the relevant inductive types, structures, and functions.
  \item State and prove theorems using tactics (\texttt{simp},
        \texttt{decide}, \texttt{omega}, \texttt{cases}).
  \item Add an \texttt{import} line to
        \texttt{specs/proofs/OmniLaTeXProofs.lean}.
  \item Run \texttt{lake build} to verify compilation and correctness.
  \item Add the module name to the verification loop in
        \texttt{lean4-ci.yml} if needed.
\end{enumerate}

Proofs should avoid \texttt{sorry} and should prefer \texttt{decide} for
finite domains and \texttt{omega} for arithmetic.  For properties that
require non-linear reasoning, consider whether the statement can be
reformulated to stay within Lean~4's \texttt{Init} library (as done in
\texttt{PageGeometry.lean}, where Float was replaced with Int for
\texttt{omega} compatibility).
